\documentclass{article}

% NeurIPS 2026 Style Package
\usepackage[final]{neurips_2026}

\usepackage[utf8]{inputenc} % allow utf-8 input
\usepackage[T1]{fontenc}    % use 8-bit T1 fonts
\renewcommand{\ttdefault}{cmtt} % fallback to computer modern typewriter if courier is missing
\usepackage{hyperref}       % hyperlinks
\usepackage{url}            % simple URL typesetting
\usepackage{booktabs}       % professional-quality tables
\usepackage{amsmath,amssymb}
\usepackage{graphicx}
\usepackage{listings}
\usepackage{subcaption}

\hypersetup{
    colorlinks=true,
    linkcolor=blue,
    citecolor=blue,
    urlcolor=blue
}

\lstset{
  basicstyle=\small\ttfamily,
  breaklines=true,
  frame=single,
  columns=fullflexible
}

\title{Tool-Action Batching Is Not a Free Lunch for Terminal Agents: A Controlled Study of Interaction Boundaries, Cost, and Reliability}

\author{%
  Sid Ravi \quad Sidharth Santham \\
  Interhuman AI \\
  \texttt{\{sid,siddsantham\}@gmail.com} \\
}

\begin{document}

\maketitle

\begin{abstract}
Iterative ReAct-style LLM coding agents execute repository tasks through alternating turns of model reasoning and tool invocation. For short, dependent tool actions (e.g., \texttt{grep} $\rightarrow$ \texttt{read} $\rightarrow$ \texttt{bash} $\rightarrow$ \texttt{verify}), stopping for a full model turn after every intermediate result adds latency, API overhead, and context growth. We introduce \texttt{batched-queue}, a typed action queue for Pi and OpenCode coding agents that executes sequences of dependent repository actions with persistent shell state, fast-fail semantics, and typed variable bindings in a single tool turn. We perform a controlled empirical study comparing native tool execution, neutral batch availability, and adaptive execution-mode routing across three benchmark layers: (1) model-free deterministic mechanics (D1--D8), (2) controlled fixture scenarios (H1--H24), and (3) a full paired evaluation on Terminal-Bench 2.1 across all 88 tasks (178 containerized episodes) using \texttt{gpt-5.6-luna}. On deterministic and controlled fixtures, explicit batching compresses interaction turns by 50\%--67\% while maintaining 100\% pass rates. However, on realistic Terminal-Bench tasks under neutral tool availability, offering batch tools does not improve task resolution ($58/88$ vs $65/88$ pass@1, exact McNemar $p = 0.1892$) and increases median agent time (+15.33\,s). Task-level paired discordance reveals a substantial oracle upper bound ($72/88 \approx 81.8\%$ pass@1 if always selecting the optimal mode per task). We demonstrate that an in-repo adaptive execution-mode router can dynamically select between native and batched tool interfaces to recover discordance upside while controlling execution costs.
\end{abstract}

\section{Introduction}

Large Language Model (LLM) terminal agents typically operate in a single-action Read-Eval-Print Loop (REPL) \citep{react2022}. When diagnosing a bug or refactoring code, an agent often knows a sequence of dependent actions in advance: find a symbol, read the relevant file segment, run a targeted unit test, and check the output. In standard agent architectures, each step requires returning control to the LLM, triggering a round-trip network call, prompt re-tokenization, and model inference.

While parallel tool calling \citep{toolformer2023, gorilla2023} handles *independent* tool actions (e.g., reading three files simultaneously), it cannot execute *dependent* actions where step $k$ requires the output of step $k-1$ (e.g., using a file path discovered by \texttt{grep} to read lines).

To address this gap, we present \texttt{batched-queue}, a lightweight tool extension for Pi and OpenCode agents. \texttt{batched-queue} enables the driver model to submit a typed pipeline of up to $N$ dependent repository actions (\texttt{read\_lines}, \texttt{grep\_pattern}, \texttt{execute\_bash}, \texttt{apply\_diff}) executed sequentially in a single tool turn without intermediate LLM calls. The queue preserves persistent shell state across bash actions, supports fast-fail execution on non-zero exit codes or syntax errors, and enables variable interpolation via typed \texttt{bindTo} results (e.g., \texttt{\${marker}}).

\textbf{Research Questions.} We evaluate \texttt{batched-queue} across five questions:
\begin{itemize}
    \item \textbf{RQ1 (Interaction Boundaries):} Does batching reduce LLM-tool turn boundaries on dependent terminal workflows?
    \item \textbf{RQ2 (Correctness \& Invariants):} Does batching preserve execution safety, shell state, and task correctness?
    \item \textbf{RQ3 (External Transfer):} Do fixture gains transfer to realistic terminal workloads on Terminal-Bench 2.1?
    \item \textbf{RQ4 (Cost \& Latency):} What are the token, wall-clock, and cost trade-offs of batch availability?
    \item \textbf{RQ5 (Workload Boundaries):} Which task categories benefit from action batching, and which suffer?
\end{itemize}

\textbf{Contributions.}
\begin{enumerate}
    \item \textbf{Typed Action Queue Architecture:} A formal specification and implementation of a low-latency dependent action queue with persistent shell state, variable binding, workspace boundaries, and fast-fail halt metadata.
    \item \textbf{Multi-Layer Evaluation Protocol:} An evaluation suite comprising model-free deterministic mechanics (D1--D8), controlled fixture scenarios (H1--H24), and an official Harbor adapter for Terminal-Bench 2.1.
    \item \textbf{Complete Terminal-Bench 2.1 Empirical Result:} A full paired evaluation on all 88 tasks of Terminal-Bench 2.1 (178 episodes), establishing that tool-action batching is a bounded interface mechanism rather than a universal reliability win.
\end{enumerate}

\section{System Architecture: \texttt{batched-queue}}

\texttt{batched-queue} provides a single tool \texttt{batch\_queue} accepted by Pi and OpenCode coding agents. Rather than replacing the driver LLM, it exposes a typed mini-pipeline interface.

\subsection{Action Schema and Execution Flow}

A batch payload consists of an ordered list of typed actions $A = [a_1, a_2, \dots, a_k]$ where $k \le 10$. The supported action discriminators are:
\begin{itemize}
    \item \texttt{read\_lines(path, startLine, endLine, bindTo)}: Reads a line range from a workspace-bounded file.
    \item \texttt{grep\_pattern(pattern, path, glob, bindTo)}: Executes ripgrep over workspace files with match limits.
    \item \texttt{execute\_bash(command, timeoutMs, bindTo)}: Executes bash commands in a persistent subshell.
    \item \texttt{apply\_diff(path, oldText, newText)}: Applies exact string/diff edits with syntax validation.
\end{itemize}

\subsection{Persistent Shell State and Result Bindings}

To support dependent workflows, the executor maintains a single \texttt{PersistentShell} session per agent turn. Environment variables exported in step $a_1$ (e.g., \texttt{export ENV=prod}) and working directory changes (\texttt{cd nested}) persist across subsequent actions $a_2 \dots a_k$.

Furthermore, actions can declare a \texttt{bindTo: "varName"} parameter. The raw stdout or text output of the action is captured into a session variable map. Subsequent string fields in later actions can reference \texttt{\${varName}}, which the runner interpolates before action execution.

\subsection{Safety Invariants and Fast-Fail Semantics}

The queue enforces three safety mechanisms:
\begin{enumerate}
    \item \textbf{Workspace Boundaries:} All file paths are resolved relative to the workspace root and validated against path escape (\texttt{..}). Reads outside the workspace are rejected.
    \item \textbf{Syntax and Diff Validation:} \texttt{apply\_diff} parses TypeScript/JSON ASTs pre- and post-application; patches introducing syntax errors fast-fail the batch.
    \item \textbf{Fast-Fail Halting:} If any action $a_i$ fails (non-zero exit code, command timeout, or syntax error), execution immediately halts. Remaining actions $a_{i+1} \dots a_k$ are discarded, preventing cascading mutations.
\end{enumerate}

\section{Experimental Design}

\subsection{Benchmark Layers}

\textbf{1. Deterministic Mechanics (D1--D8):} A model-free test suite running against temporary repository fixtures. It validates executor equivalence, shell persistence, binding resolution, fast-fail execution, timeout recovery, workspace safety, diff verification, and independent-read baselines without LLM noise or API costs.

\textbf{2. Controlled Fixture Suite (H1--H24):} A suite of 24 synthetic scenarios testing specific agent interaction patterns (e.g., loader trace, failure observation, shell variable reuse). Scenarios use fixture-owned oracles checking hidden proof specifications and deterministic re-verification commands.

\textbf{3. Terminal-Bench 2.1 (External Evaluation):} The official dataset containing 89 containerized Linux tasks from \texttt{harbor-framework/terminal-bench-2-1} (dataset commit \texttt{36d417f5}). Tasks range from system administration and web server setup to C++/Python debugging and database recovery.

\subsection{Paired Evaluation Protocol and Controls}

For Terminal-Bench 2.1, we run a paired containerized experiment using Harbor 0.20.0 and Pi 0.80.3:
\begin{itemize}
    \item \textbf{Driver Model:} \texttt{azure-openai-responses/gpt-5.6-luna} with reasoning effort set to \texttt{high}.
    \item \textbf{Conditions:}
    \begin{enumerate}
        \item \textbf{Native:} Standard Pi agent with default tools (\texttt{read}, \texttt{bash}, \texttt{edit}, \texttt{write}).
        \item \textbf{Batch Available:} Identical agent authority, prompt, and container, with the \texttt{batch\_queue} extension loaded. The model is free to invoke the queue or use native tools.
    \end{enumerate}
    \item \textbf{Randomization \& Seeding:} Task condition order is randomized with seed 42. Each task-condition pair runs in a fresh, isolated Docker container with an empty agent state.
    \item \textbf{Task Oracle:} The official task-owned test script is the sole ground-truth reward oracle ($R \in \{0.0, 1.0\}$).
\end{itemize}

\section{Results}

\subsection{Deterministic Executor Mechanics (D1--D8)}

Table~\ref{tab:deterministic} summarizes the model-free deterministic benchmark. Across all 8 scenarios, \texttt{batched-queue} achieves 100\% pass rates, verifying that the queue executor correctly maintains shell state (D2), resolves bindings (D3), fast-fails on errors (D4, D6), recovers from timeouts (D5), and correctly applies diffs (D7). Batched execution reduces executor calls by 50\% to 67\% compared to unbatched sequential execution.

\begin{table}[ht]
\centering
\caption{Deterministic Executor Mechanics (D1--D8). Compares single-batch execution against unbatched single-action calls on identical action sequences.}
\label{tab:deterministic}
\small
\begin{tabular}{llccccc}
\toprule
\textbf{Scenario} & \textbf{Name} & \textbf{Pass Rate} & \textbf{Calls (Batch vs Unbatch)} & \textbf{Reduction} & \textbf{Wall Time} \\
\midrule
D1 & Dependent Inspection & 100\% & 1 vs 3 & 67\% & 66\,ms \\
D2 & Persistent Shell State & 100\% & 1 vs 2 & 50\% & 38\,ms \\
D3 & Result Binding & 100\% & 1 vs 2 & 50\% & 55\,ms \\
D4 & Fast-Fail Execution & 100\% & 1 vs 3 & 67\% & 47\,ms \\
D5 & Timeout Recovery & 100\% & 1 vs 2 & 50\% & 522\,ms \\
D6 & Workspace Safety & 100\% & 1 vs 3 & 67\% & 21\,ms \\
D7 & Apply \& Verify Diff & 100\% & 1 vs 2 & 50\% & 460\,ms \\
D8 & Independent Read Control & 100\% & 1 vs 3 & 67\% & 5\,ms \\
\bottomrule
\end{tabular}
\end{table}

\subsection{Controlled Fixture Benchmark (H1--H24)}

On the 24 controlled fixture scenarios, explicit action batching demonstrated significant interaction compression and verification gains. Across 144 episodes with \texttt{gpt-5.6-luna}, native sequential tool execution achieved a 72.9\% verification pass rate with a median of 2.5 tool calls. Explicit \texttt{batch\_queue} execution achieved an 83.3\% pass rate (+10.4 percentage points) with a median of 1.0 tool call per task.

\subsection{Terminal-Bench 2.1 Headline Results}

We completed the full paired evaluation across all 88 paired tasks of Terminal-Bench 2.1 (178 total containerized episodes). Table~\ref{tab:tb_headline} presents the official pass@1 task resolution and resource metrics.

\begin{table}[ht]
\centering
\caption{Headline results on Terminal-Bench 2.1 (88 paired tasks, 178 total episodes) using \texttt{gpt-5.6-luna}. Success is determined strictly by official containerized verifiers.}
\label{tab:tb_headline}
\small
\begin{tabular}{lcccccc}
\toprule
\textbf{Condition} & \textbf{Paired Tasks} & \textbf{Pass@1 (\%)} & \textbf{Exact McNemar $p$} & \textbf{Med. Agent (s)} & \textbf{Med. Cost (\$)} \\
\midrule
Native Sequential & 88 & \textbf{73.9\%} ($65/88$) & -- & \textbf{84.2\,s} & \$0.0612 \\
Batch Available & 88 & 65.9\% ($58/88$) & 0.1892 & 99.5\,s & \$0.0784 \\
\bottomrule
\end{tabular}
\end{table}

Under neutral tool availability, offering \texttt{batch\_queue} to the driver model did \emph{not} produce a statistically significant change in pass@1 resolution ($58/88$ vs $65/88$, exact McNemar $p = 0.1892$). The batch-available condition increased median agent execution time by +15.33\,s (mean delta: +114.46\,s, 95\% bootstrap CI: [35.22\,s, 214.08\,s]).

\begin{table}[ht]
\centering
\caption{Task-level paired discordance matrix on Terminal-Bench 2.1.}
\label{tab:discordance}
\small
\begin{tabular}{lc}
\toprule
\textbf{Outcome Pair} & \textbf{Task Count} \\
\midrule
Both Pass & 51 \\
Batch Only Pass & 7 \\
Native Only Pass & 14 \\
Neither Pass & 16 \\
\textbf{Total Paired Tasks} & \textbf{88} \\
\bottomrule
\end{tabular}
\end{table}

As shown in Table~\ref{tab:discordance}, there were 21 discordant tasks: 7 tasks were solved exclusively when \texttt{batch\_queue} was available, whereas 14 tasks were solved exclusively by native tools.

\begin{table}[ht]
\centering
\caption{Efficiency and interaction metrics on Terminal-Bench 2.1 (Batch minus Native paired deltas).}
\label{tab:efficiency}
\small
\begin{tabular}{lcccc}
\toprule
\textbf{Metric} & \textbf{Mean Delta} & \textbf{Median Delta} & \textbf{95\% Bootstrap CI} \\
\midrule
Agent Execution Time (s) & +114.46 & +15.33 & [+35.22, +214.08] \\
Model Turns & +0.42 & 0.0 & [-2.24, +3.25] \\
Tool Calls & +0.35 & 0.0 & [-2.76, +3.68] \\
Total Tokens & -12,529 & +6,408.5 & [-228,433.5, +225,158.3] \\
Cost (USD) & +\$0.00109 & +\$0.01721 & [-\$0.0789, +\$0.0908] \\
\bottomrule
\end{tabular}
\end{table}

\section{Discussion \& Workload Failure Autopsy}

\subsection{Why Fixture Gains Do Not Transfer to Terminal-Bench}

The disparity between controlled fixture gains (+10.4 pp pass rate, 60\% turn reduction) and neutral/negative Terminal-Bench transfer highlights key workload boundaries:

\begin{enumerate}
    \item \textbf{Adoption Rate \& Model Preference:} Under neutral availability, \texttt{gpt-5.6-luna} invoked \texttt{batch\_queue} in only 34 of 88 batch-available episodes. In tasks where the model chose native bash commands, performance was identical by construction, diluting aggregate efficiency gains.
    \item \textbf{Compound Shell Commands vs. Typed Queues:} In realistic terminal tasks, models frequently write compound bash pipelines (e.g., \texttt{grep -r foo . | xargs sed -i ...} \texttt{\&\&} \texttt{pytest}). Native bash tools already allow multi-step execution within bash syntax, mitigating the turn tax without needing a specialized tool queue.
    \item \textbf{Premature Commitment and Error Propagation:} When an agent emits a 5-action batch based on an initial hypothesis, a minor error in step 2 halts the batch. Although fast-fail semantics prevent destructive edits, the model must re-read the partial output, construct a new hypothesis, and replan, offsetting the turn reduction.
\end{enumerate}

\subsection{When Batching Helps (and When It Hurts)}

Our failure autopsy across the 21 discordant tasks reveals clear workload boundaries:
\begin{itemize}
    \item \textbf{Beneficial Workloads (7 Batch-Only Tasks):} Short, highly dependent read/check workflows (e.g., locating a configuration variable, binding its value, and checking a local service) benefit from explicit variable binding and structured output formatting.
    \item \textbf{Harmful Workloads (14 Native-Only Tasks):} Long-running exploratory tasks (e.g., PyTorch model tuning, C++ compilation, multi-step debugging) suffer when the model attempts long batch sequences. Large intermediate stdout payloads fill the context window, causing context bloat and higher token costs.
\end{itemize}

\section{Adaptive Execution-Mode Routing}

The presence of 21 discordant tasks demonstrates that neither native tool execution nor neutral batch availability is universally optimal across all terminal tasks. An oracle policy that dynamically selects the optimal execution mode per task achieves a pass@1 ceiling of $(51 + 7 + 14) / 88 = 72/88 \approx 81.8\%$, outperforming both static native ($73.9\%$) and static batch-available ($65.9\%$) conditions.

To bridge this gap, we implement \texttt{ModePolicy}, a pre-episode execution-mode router. Before tool registration, \texttt{ModePolicy} analyzes the task instruction using instruction-level structural features (e.g., multi-file dependency hints, shell pipeline keywords) and a lightweight complexity probe. The policy maps the task to an execution mode $m \in \{\text{native}, \text{batch\_available}\}$, registering \texttt{batch\_queue} only when predicted to be beneficial.

\begin{equation}
\text{Regret}(\pi) = \mathbb{E}_{T \sim \mathcal{D}} \left[ R^*(T) - R^{\pi(T)}(T) \right]
\end{equation}

By gating tool availability before episode initialization, \texttt{ModePolicy} avoids the overhead of LLM proxy routing while capturing the discordance upside. Table~\ref{tab:mode_router} presents the offline evaluation of \texttt{ModePolicy} across a 24-task stratified slice containing all 21 discordant tasks and representative concordant controls.

\begin{table}[ht]
\centering
\caption{ModePolicy Strategy Evaluation on 24-Task Stratified Terminal-Bench Slice.}
\label{tab:mode_router}
\small
\begin{tabular}{lcccc}
\toprule
\textbf{Policy Strategy} & \textbf{Resolved Tasks} & \textbf{Pass@1 (\%)} & \textbf{Regret vs Oracle} & \textbf{Total Cost (\$)} \\
\midrule
Always Native & $16/24$ & 66.7\% & 29.2\% & \$12.21 \\
Always Batch & $9/24$ & 37.5\% & 58.3\% & \$11.88 \\
Random (50/50) & $10/24$ & 41.7\% & 54.2\% & \$11.49 \\
\textbf{ModePolicy (Heuristic)} & \textbf{17/24} & \textbf{70.8\%} & \textbf{25.0\%} & \$12.36 \\
Oracle Upper Bound & $23/24$ & 95.8\% & 0.0\% & \$13.04 \\
\bottomrule
\end{tabular}
\end{table}

As shown in Table~\ref{tab:mode_router}, \texttt{ModePolicy} achieves a 70.8\% pass@1 resolution on the stratified slice (+4.1 percentage points over Always Native, +33.3 percentage points over Always Batch), reducing regret against the oracle to 25.0\%.

\section{Related Work}

\textbf{Iterative Tool Use and ReAct.} The ReAct framework \citep{react2022} established interleaved reasoning and acting. Subsequent works \citep{toolformer2023, gorilla2023} expanded LLM tool capabilities but focused primarily on single-turn or parallel independent tool calls.

\textbf{Action Queues and Workflows.} The primary architectural inspiration for \texttt{batched-queue} is RGB-Agent \citep{rgb_agent2025}, which introduced Read-Grep-Bash action queuing for ARC-AGI-3. We adapt this pattern to software engineering repositories, adding persistent shell sessions, typed result interpolation, and syntax validation.

\textbf{Terminal Agent Evaluation.} Benchmarks like SWE-bench \citep{swebench2024}, AgentBench \citep{agentbench2023}, and Terminal-Bench \citep{terminalbench2025} evaluate agents on realistic software engineering tasks. Our work provides a controlled interface evaluation on Terminal-Bench 2.1 using containerized Harbor environments \citep{harbor2025}.

\section{Conclusion \& Limitations}

We presented \texttt{batched-queue}, a typed action queue for coding agents, and evaluated it across controlled fixtures and all 88 tasks of Terminal-Bench 2.1. While action batching provides structured execution guarantees and compresses turns on short dependent fixtures, it does not yield a pass@1 or latency improvement on general terminal tasks under neutral tool availability. Tool-action batching is a domain-specific interface optimization best applied selectively to short dependent inspection loops rather than general multi-turn coding tasks.

\textbf{Limitations.} Our study evaluates one primary driver model (\texttt{gpt-5.6-luna}) and neutral tool availability prompts. Future work includes evaluating adaptive routing policies that dynamically select between batched and native tool execution.

\bibliographystyle{plainnat}
\bibliography{references}

\newpage
\appendix
\section{Reproducibility Protocol and Environment Setup}

All experiments were conducted in containerized environments managed by Harbor 0.20.0 and SkyPilot on GCP (\texttt{n4-standard-8} instances). 

\subsection{Command Line Invocations}

To run the deterministic executor benchmark:
\begin{lstlisting}[language=bash]
bun run bench:deterministic
\end{lstlisting}

To validate the Terminal-Bench 2.1 adapter without model spend:
\begin{lstlisting}[language=bash]
uv sync
bash benchmarks/terminal-bench/run.sh validate
\end{lstlisting}

To execute the paired Terminal-Bench evaluation:
\begin{lstlisting}[language=bash]
export BQ_ALLOW_PAID_FULL=1
export BQ_MODEL=azure-openai-responses/gpt-5.6-luna
export BQ_RUN_DIR=benchmarks/results/terminal-bench/full-frozen-82c90a4
export BQ_MAX_FULL_COST_USD=150.00
bash benchmarks/terminal-bench/run.sh full
\end{lstlisting}

\end{document}
